\documentclass[paper=a4, fontsize=11pt]{article}

%----------------------------------------------------------------------------------------
%	PACKAGES AND OTHER DOCUMENT CONFIGURATIONS
%----------------------------------------------------------------------------------------

\usepackage{amsmath,amsfonts,amsthm} % Math packages
\usepackage{sectsty} % Allows customizing section commands
\allsectionsfont{\centering \normalfont\scshape} % Make all sections centered, the default font and small caps

\usepackage{fancyhdr} % Custom headers and footers
\pagestyle{fancyplain} % Makes all pages in the document conform to the custom headers and footers
\fancyhead{} % No page header - if you want one, create it in the same way as the footers below
\fancyfoot[L]{} % Empty left footer
\fancyfoot[C]{} % Empty center footer
\fancyfoot[R]{\thepage} % Page numbering for right footer
\renewcommand{\headrulewidth}{0pt} % Remove header underlines
\renewcommand{\footrulewidth}{0pt} % Remove footer underlines
\setlength{\headheight}{13.6pt} % Customize the height of the header

\numberwithin{equation}{section} % Number equations within sections (i.e. 1.1, 1.2, 2.1, 2.2 instead of 1, 2, 3, 4)
\numberwithin{figure}{section} % Number figures within sections (i.e. 1.1, 1.2, 2.1, 2.2 instead of 1, 2, 3, 4)
\numberwithin{table}{section} % Number tables within sections (i.e. 1.1, 1.2, 2.1, 2.2 instead of 1, 2, 3, 4)


\usepackage{graphicx}
\graphicspath{ {images/} }

\usepackage{xepersian}
\settextfont[Path=fonts/]{Vazir.ttf}
%\setlatintextfont{Times New Roman}

%----------------------------------------------------------------------------------------
%	TITLE SECTION
%----------------------------------------------------------------------------------------

\newcommand{\horrule}[1]{\rule{\linewidth}{#1}} % Create horizontal rule command with 1 argument of height

\title{
\normalfont\normalsize
\includegraphics[scale=0.1]{aut}
\hspace{5cm}
\includegraphics[scale=0.1]{ceit} \\
\textsc دانشگاه صنعتی امیرکبیر \\
\textsc دانشکده مهندسی کامپیوتر و فناوری اطلاعات
\horrule{0.5pt} \\ [0.4cm] % Thin top horizontal rule
\huge سیستم‌های توزیعی \\ % The assignment title
\huge فاز اول پروژه دوم \\ % The assignment title
\horrule{2pt} \\ [0.5cm] % Thick bottom horizontal rule
}

\author{پرهام الوانی}

\date{\normalsize\today} % Today's date or a custom date

\begin{document}

\maketitle % Print the title

\paragraph{}
برای پیاده‌سازی پروژه سه گام در نظر گرفته شده است. در اولین گام ارتباط بین سرور و کلاینت پیاده‌سازی شد
و در گام ارتباط بدون خطا فرض می‌شود.
در دومین و سومین گام به ترتیب ارتباط دارای خطا و بسته شدن صحیح سرور و کلاینت پیاده‌سازی گشت.

\paragraph{}
دو قسمت اصلی در کلاینت در نظر گرفته می‌شود، قسمت اول \lr{receiver}
است که داده‌ها را از سوکت دریافت کرده و در کانال \lr{incoming}
می‌ریزد. قسمت دوم \lr{handler}
نام دارد که در واقع قلب سیستم است و تمام کنترل‌ها در آنجا صورت می‌گیرد
دلیل این امر جلوگیری از وقوع \lr{race condition}
می‌باشد.

\paragraph{}
\lr{handler}
داده‌ها را از کانال \lr{incoming} دریافت می‌کند
و آن‌ها را در یک \lr{hash map} بر اساس شماره دنباله‌شان ذخیره می‌کند.
در ادامه پیام‌های \lr{ack} مناسب را تولید می‌کند.
پیام‌هایی که ترتیب آن‌ها درست است در کانال \lr{rmsg} برای دریافت توسط لایه‌ی بالاتر قرار می‌گیرند.

لایه اپلیکشن هم پیام‌هایی که می‌خواهد ارسال شوند را در کانال \lr{tmsg} قرار می‌دهد.
\lr{handler} این پیام‌ها را نیز در یک \lr{hash map} بر اساس شماره دنباله‌شان قرار می‌دهد
با این تفاوت که طول این \lr{hash map} هرگز نباید از \lr{window size (windows)}
بزرگتر شود.
با دریافت هر \lr{ack} آن پیام از لیست پیام‌های در حال ارسال حذف می‌شود.

\lr{handler}
با رویدادهای زمانی که در بازه‌های زمانی \lr{epochMillis} رخ می‌دهند
سه کار انجام می‌دهد. در صورتی که هیچ پیامی دریافت نشده باشد \lr{ack}
با شماره دنباله‌ی صفر ارسال می‌کند و اگر داده‌هایی در حال ارسال داشته باشد
تمام آن‌ها را ارسال می‌کند. در صورتی که هنوز ارتباط با سرور شکل نگرفته باشد پیام
\lr{connect}
را باز تکرار می‌کند.

برای بستن ارتباط یک کلاینت دو حالت وجود دارد. در اولین حالت
ارتباط با سرور \lr{timeout} می‌کند که در این حالت
ارتباط بدون توجه به وضعیت پیام‌ها بسته می‌شود و پیام خطا در
کانال \lr{err} قرار می‌گیرد.
در دومین حالت ارتباط به درخواست کاربر بسته می‌شود که در این حالت
نیاز است که \lr{handler}
تمام پیام‌هایی که در صف بافی مانده است را ارسال کند.
برای مدیریت این شرایط از متغیر \lr{status} استفاده می‌کنیم
\lr{status}
در واقع یک مانیتور است که دسترسی به یک \lr{integer}
را مدیریت می‌کند.

\end{document}
